\scanpage{78}
\mspage{74}
Cela montre que pour la construction

2. \uline{Spatios} Deux façons de \unsure{le|les} décrire.

\deleted{\unsure{Pr|P}}\uline{Définition 1} Un spatios est un ensemble $\Sigma$, muni \unsure{de$'$|de} \deleted{deux} relation\deleted{s} $\leq$, $\deleted{\unsure{\mathrel{|\circ|}|\mathrel{\#}}}$ et d'un endomorphisme $C$ (relation aussi, si on veut !), satisfaisant les conditions

\uline{Spat 1} La relation $\leq$ est une \uline{relation d'ordre}, admettant des $\operatorname{Sup}$ et des $\operatorname{Inf}$ quelconques (donc un plus petit élément, noté \inserted{$0$ ou} $0_{\Sigma}$, et un plus grand élément, noté \inserted{$1$ ou} $1_{\Sigma}$).

\uline{Spat 2} $C$ est une \uline{involution} ($C\deleted{\unsure{x}}(Cx)=x$) et \deleted{et} ``renverse la relation d'ordre'' :
\[
\deleted{\unsure{C}}x\leq y\Longleftrightarrow Cy\leq Cx.
\]

\deleted{\uline{\unsure{Spat 3}} $\forall x$}
\[
\left[\begin{aligned}
&\text{\uline{Cor} 1}\quad C(\operatorname{Sup}_{i}x_{i})=\operatorname{Inf}_{i}Cx_{i},\quad C(\operatorname{Inf}_{i}x_{i})=\operatorname{Sup}_{i}Cx_{i}.\\
&\text{\uline{Cor. 2}}\quad C0=1,\quad C1=0
\end{aligned}\right.
\]

\deleted{A un Spatios est \unsure{associé|associée} une structure de disjonction canonique, en \unsure{posant}}
\[
\deleted{x\mathrel{|\unsure{\circ|0}|}y\mathrel{\overset{\mathrm{def}}{\Longleftrightarrow}}x\leq Cy\quad(\Longleftrightarrow y\leq Cx\ \text{par Spat 2}).}
\]

\uline{Spat 3} $\forall x\in\Sigma$, on a
\[
x\wedge Cx=0_{\deleted{\unsure{E}}\Sigma}
\]
\unsure{ce|ou} qui \unsure{équiv|équiv.} par Spat 2 (Cor. 1 et 2) :
\[
x\unsure{:}\vee Cx=1_{\Sigma}.
\]
% archive sheet number: 77
